\section{The model and the variational formula}\label{sec:model}

Fix a pattern density $\alpha>0$. Let $\nu_N$ be the uniform \emph{probability}
measure on $\{-1,1\}^N$, let $M_N=\lfloor\alpha N\rfloor$, and let
$g^1,\ldots,g^{M_N}$ be independent standard Gaussian vectors in
$\mathbb R^N$. For a bounded Borel function $f:\mathbb R\to\mathbb R$, put
\begin{equation}\label{def:pressure}
 b_a(x)=\frac{g^a\cdot x}{\sqrt N},\qquad
 Z_N(f)=\int\exp\!\left\{\sum_{a=1}^{M_N}f(b_a(x))\right\}\nu_N(dx),
 \qquad p_N(f)=\frac1N\log Z_N(f).
\end{equation}
The usual inverse temperature and activation are recovered by writing
$f=\beta\phi$, with $\beta>0$. Counting measure instead of $\nu_N$ adds
$\log2$ to the pressure. With either convention, physical free energy per
spin is the corresponding pressure multiplied by $-1/\beta$.

\subsection{Two finite Gaussian recursions}

An overlap path is a nondecreasing function $q:(0,1)\to[0,1]$. A trial field
path is a nonnegative nondecreasing step function $h:(0,1)\to[0,\infty)$
with finite values. Paths are identified up to equality almost everywhere;
all path integrals and $L^1$ norms use Lebesgue measure. Write $\mathcal Q$
for the overlap paths and $\mathcal H_{\mathrm{step}}$ for the trial fields.

A step path is described on a partition
\begin{equation}\label{def:partition}
 0=\zeta_0<\zeta_1<\cdots<\zeta_k=1,
 \qquad w_i=\zeta_{i+1}-\zeta_i,\quad 0\le i<k.
\end{equation}
Its value on $(\zeta_i,\zeta_{i+1})$ carries subscript $i$; repeated path
values are allowed. For a standard normal random variable $G$, define
\begin{equation}\label{def:transform}
 (\mathcal T_{s,d}U)(x)=
 \begin{cases}
 d^{-1}\log\mathbb E\exp\{dU(x+\sqrt s\,G)\},&d>0,\\
 \mathbb E U(x+\sqrt s\,G),&d=0,
 \end{cases}
 \qquad s\ge0,\quad 0\le d\le1.
\end{equation}

\begin{definition}[Pattern and field recursions]\label{def:recursions}
For a step overlap path $q$, append $q_{-1}=0$ and $q_k=1$. Starting at
$U_k=f$, set
\begin{equation}\label{def:pattern-recursion}
 U_{i-1}=\mathcal T_{q_i-q_{i-1},\zeta_i}U_i
 \quad(i=k,k-1,\ldots,0),\qquad V_f(q)=U_{-1}(0).
\end{equation}
For a step field path $h$, append $h_{-1}=0$. Starting instead at
$U_{k-1}(x)=\log\cosh x-h_{k-1}/2$, set
\begin{equation}\label{def:field-recursion}
 U_{i-1}=\mathcal T_{h_i-h_{i-1},\zeta_i}U_i
 \quad(i=k-1,k-2,\ldots,0),\qquad \ell(h)=U_{-1}(0).
\end{equation}
The letter $U$ refers to the recursion under discussion.
\end{definition}

The initial Gaussian average has coefficient $\zeta_0=0$. The pattern
recursion also has a final Gaussian average with coefficient $\zeta_k=1$
and variance $1-q_{k-1}$; the field recursion has no such extra interval.
For example, for constant paths $q(u)=r$ and $h(u)=H$,
\begin{align}
 V_f(r)&=\mathbb E_G\log\mathbb E_{G'}
   \exp f(\sqrt r\,G+\sqrt{1-r}\,G'),\label{def:one-level-pattern}\\
 \ell(H)&=\mathbb E\log\cosh(\sqrt H\,G)-H/2,
 \label{def:one-level-field}
\end{align}
where $G,G'$ are independent standard normals.

Subdividing an interval without changing its path value inserts a
zero-variance transform, so the definitions do not depend on the chosen
partition. For $f\in C_c^\infty(\mathbb R)$, the pattern functional is Lipschitz in the
$L^1$ distance between step paths. For bounded Borel $f$, it has a unique
continuous extension to $\mathcal Q$, constructed by uniform approximation
of the pattern functionals in Lemma~\ref{ext:path-approx}. We use that extension throughout. The field
functional is $1/2$-Lipschitz on step fields and therefore extends to bounded
nonnegative nondecreasing fields; see Proposition~\ref{var:continuity-order}.

\begin{definition}[Ising entropy]\label{def:entropy}
For $q\in\mathcal Q$, define the extended-valued functional
\begin{equation}\label{def:entropy-formula}
 S_{\mathrm I}(q)=\sup_{h\in\mathcal H_{\mathrm{step}}}
 \left\{\ell(h)+\frac12\int_0^1h(u)q(u)\,du\right\}.
\end{equation}
The field partition is independent of the overlap partition. In calculations
with two step paths, we always pass to a common refinement.
\end{definition}

\begin{theorem}[Gaussian Ising perceptron]\label{thm:main}
For every $\alpha>0$ and every bounded Borel $f:\mathbb R\to\mathbb R$,
\begin{equation}\label{thm:formula}
 \lim_{N\to\infty}\mathbb E p_N(f)
 =\mathcal P_f(\alpha)
 :=\inf_{q\in\mathcal Q}\{\alpha V_f(q)+S_{\mathrm I}(q)\}.
\end{equation}
Moreover, $p_N(f)$ converges to $\mathcal P_f(\alpha)$ in probability.
The variational value is finite. In particular, the conclusion holds for
$f=\beta\phi$ at every $\beta>0$ and every bounded continuous activation
$\phi$, without symmetry, concavity, or a smallness restriction.
\end{theorem}

The elementary bounds $S_{\mathrm I}\ge0$, $S_{\mathrm I}(0)=0$ and
$|V_f|\le\|f\|_\infty$ will be proved below. Since
$V_f(0)=\log\mathbb E e^{f(G)}$, they give
\begin{equation}\label{def:finite-value}
 -\alpha\|f\|_\infty\le\mathcal P_f(\alpha)
 \le\alpha\log\mathbb E e^{f(G)}.
\end{equation}
The proof of Theorem~\ref{thm:main} first assumes $f\in C_c^\infty(\mathbb R)$.
The upper bound is Proposition~\ref{upper:bound}, and
Proposition~\ref{cav:lower-bound} gives the matching lower bound.
Section~\ref{sec:extensions} removes smoothness
and proves the remaining convergence claims. It also gives extensions to
integrable unbounded activations and to independent random-sign patterns.

\subsection{Probabilistic conventions}

A Gibbs measure is a reference probability reweighted by an exponential
energy and then normalized. Replicas are independent samples from that
measure conditional on its disorder and reference data. Brackets
$\langle\cdot\rangle$ denote replica expectation, and $\mathbb E$ averages
the disorder and any explicitly introduced random data. A deterministic
perturbation parameter is held fixed under $\mathbb E$ unless its averaging
is specified. The spin overlap is
\[
 R(x,x')=N^{-1}x\cdot x',\qquad R_{ab}=R(x^a,x^b).
\]
An array limit means convergence in law of every finite sampled restriction,
including disorder averaging. All constants may depend on $\alpha$ and the
fixed smooth activation. Dependence on a fixed number of cavity coordinates
$L$ is allowed only where indicated; the limit in $N$ is always taken before
letting $L$ increase.
